\begin{textsample}{POS Dim 1 – 2018\_c4 – Score -3.00 – 2018\_c4/enade\_2018\_dir\_72417de982.txt}  \label{ex:f1_pos_302}
CAPÍTULO 4 PERCEPÇÃO DA PROVA Nas análises feitas neste capítulo, constam as percepções dos concluintes da Área de Direito sobre a prova aplicada no Enade/2018.
Essas percepções foram mensuradas por meio de nove questões que avaliaram, desde o grau de dificuldade da prova até o tempo gasto para concluí -la.
As percepções sobre a prova foram relacionadas ao desempenho dos estudantes e à Grande Região de funcionamento do curso.
O questionário de percepção da prova encontra -se ao final do Anexo VII, que traz a reprodução do exame.
O desempenho dos estudantes foi classificado em quatro quartos.
Para tanto, esse desempenho foi ordenado de forma ascendente.
O percentil 25, P25, também conhecido como primeiro quartil, é a nota de desempenho que deixa um quarto ( 25 \% ) dos valores observados abaixo e três quartos acima.
A Figura 1 apresenta uma ilustração desse conceito.
O quarto inferior de desempenho é composto pelas notas abaixo do primeiro quartil.
Já o percentil 75, P75, também conhecido como terceiro quartil, é o valor para o qual há três quartos ( 75 \% ) dos dados abaixo e um quarto acima dele.
O quarto superior de desempenho é composto pelas notas iguais ou acima do terceiro quartil.
O percentil 50, P50, também conhecido como mediana, é o valor que divide as notas em dois conjuntos de igual tamanho.
O segundo quarto inclui valores entre o primeiro quartil ( P25 ) e a mediana.
O terceiro quarto contém os valores entre a mediana ( P50 ) e o terceiro quartil ( P75 ).
Vale ressaltar que percentis, quartis e medianas são pontos que, não obrigatoriamente, pertencem ao conjunto original de dados, ao passo que os quartos são subconjuntos dos dados originais.
A seguir, serão apresentados gráficos com resultados selecionados, relativos às nove questões avaliadas por grupos de estudantes.
Nas barras dos gráficos, apresenta -se o percentual de alunos que assinalaram uma das opções ou a soma das porcentagens daqueles que assinalaram duas ( ou três ) delas.
Por exemplo, para as questões 1 e 2, nos gráficos, é apresentada a porcentagem total de participantes que assinalaram as opções ( D ) Difícil e ( E ) Muito difícil.
Em cada barra, foram assinalados, também, os extremos do intervalo de confiança de 95 \% como linhas verticais unidas por uma linha horizontal na forma da letra H maiúscula.
O estimador de um parâmetro com um certo nível de confiança ( e.g. 95,0 \% ) deve conter o parâmetro no intervalo de confiança em 95 \% das vezes.
Na comparação entre os estimadores dos parâmetros de duas classes de uma dada categoria ( e.g.
Norte e Nordeste nas Grandes Regiões ou de primeiro e último quarto dentro de desempenho ) associados aos seus respectivos intervalos de confiança, diz -se que não há diferença estatisticamente significativa entre os parâmetros das duas categorias quando há interseção entre os intervalos de confiança e que há diferença, se os intervalos de confiança são disjuntos ( para maiores informações vide Glossário ).
Nas Tabelas no Anexo II, são apresentados os valores absolutos e a distribuição percentual18 das alternativas válidas das nove questões, segundo o mesmo recorte de desempenho dos alunos e Grande Região de funcionamento do curso. 4.1 GRAU DE DIFICULDADE DA PROVA 4.1.1 Componente de Formação Geral Ao avaliarem Qual o grau de dificuldade desta prova na parte de Formação Geral? ( Questão 1 ), 36,1 \% do grupo de inscritos e presentes optaram pelas alternativas Difícil ou Muito difícil.
Entretanto, para mais da metade dos estudantes ( 52,4 \% ), o Componente de Formação Geral da prova foi considerado com grau de dificuldade Médio ( Gráfico 4.1, Gráfico 4.2 e, no Anexo II, a Tabela II.1 ).
O percentual de estudantes que consideraram a prova Difícil ou Muito difícil foi maior na região Centro-Oeste, onde a proporção foi de 39,4 \%, enquanto a de menor incidência foi a Nordeste, com 31,3 \%.
No Gráfico 4.1, é possível observar que a diferença entre as regiões Norte e Sudeste e a diferença entre as regiões Sul e Centro-Oeste não são estatisticamente significativas, todas as demais diferenças regionais são estatisticamente significativas.
Nas Grandes Regiões, a proporção de presentes à prova que consideraram o Componente de Formação Geral como sendo de grau de dificuldade Médio esteve entre 49,8 \%, na região Sul, e 56,4 \%, na região Nordeste.
O percentual de alunos que consideraram a prova Difícil ou Muito difícil foi decrescente em função dos quartos de desempenho : 41,9 \%, no primeiro quarto e 27,7 \%, no quarto, grupo de melhor desempenho na prova.
Nos quartos de desempenho intermediários, a proporção de alunos que consideraram a prova Difícil ou Muito difícil foi de 39,6 \%, no segundo quarto e de 35,3 \%, no terceiro quarto.
As diferenças entre todos os quartos de desempenho são estatisticamente significativas.
Para todos os quartos de desempenho, a alternativa modal para essa pergunta foi a Médio, com 48,7 \% e 55,0 \% dos respondentes nos quartos extremos, primeiro e quarto, respectivamente.
Estas proporções são crescentes com o desempenho. 4.1.2 Componente de Conhecimento Específico Ao responderem à Questão 2 – Qual o grau de dificuldade desta prova na parte de Componente Específico? – 47,9 \% do grupo de estudantes classificaram -na como Difícil ou Muito difícil.
Além disso, o Componente de Conhecimento Específico da prova foi considerado com grau de dificuldade Médio por 47,3 \% dos alunos ( Gráfico 4.3, Gráfico 4.4, e, no Anexo II, a Tabela II.2 ).
A análise das respostas dos estudantes quanto ao grau de dificuldade do Componente de Conhecimento Específico da prova, agregado por Grande Região, mostra que as diferenças entre as proporções de alunos que a avaliaram como Difícil ou Muito difícil são estatisticamente significativas entre todas regiões com exceções das diferenças entre as regiões Norte e Sudeste e Sul e Centro-Oeste.
O percentual de alunos que classificaram o grau de dificuldade como Médio, no Componente de Conhecimento Específico, variou de 43,6 \% a 52,1 \%, para as regiões Sul e Nordeste, respectivamente.
Considerando -se a avaliação da dificuldade das questões do Componente de Conhecimento Específico da prova, de acordo com o desempenho dos estudantes, observa -se que não há diferença estatisticamente significativa dos resultados dentre os quartos de desempenho extremos ( primeiro e último ) e dentre os quartos intermediários de desempenho ( segundo e terceiro ).
A proporção dos que classificaram a parte específica como Difícil ou Muito difícil variou de 46,5 \% ( quarto inferior e quarto superior, cada ) a 49,4 \% ( terceiro quarto ).
Já a proporção dos que responderam que o grau de dificuldade das questões do Componente de Conhecimento Específico da prova foi Médio variou de 46,6 \%, nos segundo e terceiro quartos a 48,7 \%, no quarto superior. 4.2 EXTENSÃO DA PROVA EM RELAÇÃO AO TEMPO TOTAL Indagados quanto à extensão da prova, em relação ao tempo total oferecido para a sua resolução ( Questão 3 ), os estudantes apontaram, com maior incidência, a alternativa que considerava a extensão adequada, para todas as agregações consideradas ( Gráfico 4.5, Gráfico 4.6, e, no Anexo II, a Tabela II.3 ).
O percentual de alunos que responderam ser a extensão da prova adequada foi de 47,6 \%.
Já 44,0 \% dos inscritos presentes consideraram que a prova foi longa ou muito longa, e 8,3 \% a avaliaram como curta ou muito curta.
Dentre as Grandes Regiões, a proporção daqueles que avaliaram a prova como longa ou muito longa, em relação ao tempo total destinado à sua resolução, variou de 41,0 \%, na região Norte, até 44,9 \%, na região Centro-Oeste.
Há diferenças estatisticamente significativas da região Norte com as demais regiões.
Considerando -se o desempenho dos alunos, nota -se, ainda, que 46,0 \% consideraram a extensão da prova adequada no quarto de desempenho inferior e 49,7 \% no de melhor desempenho ( quarto superior ).
Nos quartos intermediários, essa proporção foi de 46,8 \%, no segundo quarto e de 47,8 \%, no terceiro.
No Gráfico 4.6, pode -se constatar que há diferença estatisticamente significativa entre as proporções de estudantes que consideraram a prova longa ou muito longa em função dos quartos de desempenho, sendo a maior proporção dessa resposta no primeiro quarto de desempenho ( 45,0 \% ), e a menor, no quarto superior ( 42,1 \% ). 4.3 COMPREENSÃO DOS ENUNCIADOS DAS QUESTÕES 4.3.1 Componente de Formação Geral Em relação aos enunciados das questões do Componente de Formação Geral ( Questão 4 ), as opiniões foram positivas, já que 66,5 \% dos alunos avaliados consideraram todos ou a maioria dos enunciados das questões claros e objetivos ( Gráfico 4.7, Gráfico 4.8, e, no Anexo II, a Tabela II.4 ).
Na análise regional, a percentagem de estudantes que avaliaram que todos ou a maioria dos enunciados das questões do Componente de Formação Geral estavam claros e objetivos variou de 65,1 \%, na região Centro-Oeste, a 67,9 \%, na região Sul.
A análise das percepções dos estudantes sobre a clareza e objetividade dos enunciados permite afirmar que todos, ou a maioria dos enunciados de questões relativas ao Componente de Formação Geral, foram considerados claros e objetivos para a maior parte dos respondentes ( maior ou igual a 65,1 \% em todas as regiões e maior ou igual a 58,3 \% para todos os quartos de desempenho ).
Segundo o desempenho, observa -se que a proporção dos que emitiram essa opinião cresce conforme o desempenho aumenta, com diferenças estatisticamente significativas entre todos os quartos de desempenho.
No quarto superior, a clareza e objetividade de todos ou da maioria dos enunciados das questões foram percebidas por 75,3 \%. 4.3.2 Componente de Conhecimento Específico Em relação aos enunciados das questões do Componente de Conhecimento Específico da prova, para 65,6 \% dos estudantes avaliados da Área de Direito, a clareza e a objetividade ( Questão 5 ) estavam presentes em todas ou na maioria das questões ( Gráfico 4.9, Gráfico 4.10, e no Anexo II, a Tabela II.5 ).
A maioria dos estudantes de todas as Grandes Regiões brasileiras considerou claros e objetivos todos ou a maioria dos enunciados das questões do Componente de Conhecimento Específico da prova, percentual sempre maior ou igual a 64,4 \%.
As diferenças entre as regiões Sudeste e as demais regiões são estatisticamente significativa, exceto quando confrontada com a região Norte.
A proporção de estudantes que consideraram os enunciados das questões claros e objetivos apresenta uma tendência crescente em relação ao aumento de desempenho : mais elevada no quarto superior ( 72,8 \% ), se comparada ao quarto inferior de desempenho ( 59,3 \% ).
As diferenças entre os quartos de desempenhos são todas estatisticamente significativas. 4.4 SUFICIÊNCIA DAS INFORMAÇÕES/INSTRUÇÕES FORNECIDAS Ao avaliarem as informações/instruções fornecidas para a resolução das questões ( Questão 6 ), 78,8 \% dos respondentes da Área de Direito de todo o Brasil afirmaram que estas eram até excessivas ou suficientes em todas ou na maioria das questões ( Gráfico 4.11, Gráfico 4.12, e, no Anexo II, a Tabela II.6 ).
Quanto à distribuição de respondentes pelas Grandes Regiões, observa -se que a proporção de estudantes que consideraram as informações/instruções fornecidas até excessivas ou suficientes em todas ou na maioria das questões foi sempre superior ou igual a 77,0 \%, na região Nordeste chegando a 81,4 \%, na região Sul.
Há diferenças estatisticamente significativas entre a região Sudeste com as demais.
Levando -se em conta o desempenho dos participantes, observa -se uma tendência crescente das proporções de participantes que avaliaram as informações/instruções como até excessivas ou suficientes em todas ou na maioria das questões, havendo diferenças estatisticamente significativas entre o quarto inferior e os dois quartos superiores, como mostrado no Gráfico 4.12.
O percentual foi mais elevado no quarto superior ( 84,9 \% ), percentual superior à média nacional ( 78,8 \% ).
Já no quarto inferior, a suficiência das informações/instruções declarada como até excessiva, em todas ou na maioria das questões foi percebida por 72,8 \% dos respondentes. 4.5 DIFICULDADE ENCONTRADA AO RESPONDER À PROVA Indagados sobre as dificuldades com as quais se depararam ao responder à prova ( Questão 7 ), 13,7 \% dos estudantes apontaram o Desconhecimento do conteúdo.
Para 48,9 \%, a Forma diferente de abordagem do conteúdo foi indicada como dificuldade.
Já a Falta de motivação para fazer a prova foi a dificuldade apontada por 14,3 \% dos respondentes.
Considerando -se todo o Brasil, 16,2 \% dos respondentes afirmaram que não tiveram qualquer tipo de dificuldade para responder à prova ( Tabela II.7 no Anexo II ).
Nos Gráficos 4.13 e 4.14, são apresentados os percentuais de estudantes que apontaram o Desconhecimento do conteúdo como dificuldade ao responder à prova.
Na análise por Grandes Regiões, o percentual de inscritos e presentes que apontaram o Desconhecimento do conteúdo como dificuldade ao responder à prova não superou 15,4 \%.
Os percentuais variaram de 12,1 \%, na região Nordeste, a 15,4 \%, na região Sul.
A Forma diferente de abordagem do conteúdo foi a escolha modal dos estudantes, com percentuais que variaram de 47,2 \% ( região Nordeste ) a 51,5 \% ( região Centro-Oeste ).
O percentual de alunos que citaram a Falta de motivação como dificuldade variou de 13,1 \% ( região Norte ) a 16,5 \% ( região Nordeste ).
Os que declararam não ter qualquer tipo de dificuldade para responder à prova variaram de 14,3 \%, na região Centro-Oeste, a 16,9 \%, na Sudeste.
Em relação aos quartos de desempenho, o Desconhecimento do conteúdo foi a opção escolhida por 14,4 \% dos estudantes do quarto superior e por 13,9 \% do quarto inferior.
Não são observadas diferenças estatisticamente significativas entre as proporções dos quartos de desempenho, exceto entre os quartos de desempenhos intermediários e o quarto superior de desempenho A alternativa modal para os alunos, quando agregados pelos quartos de desempenho, foi que a dificuldade encontrada foi causada pela Forma diferente de abordagem do conteúdo : 46,7 \% no quarto inferior e 46,4 \%, do quarto superior escolheram essa alternativa.
Nos quartos intermediários, 51,2 \% dos alunos apontaram essa dificuldade, média superior a nacional. 4.6 CONTEÚDOS DAS QUESTÕES OBJETIVAS DA PROVA Ao analisar os conteúdos das questões objetivas da prova ( Questão 8 ), um percentual muito pequeno dos estudantes avaliados, apenas 4,2 \%, afirmou que não estudou ainda a maioria desses conteúdos ( Gráficos 4.15, Gráfico 4.16, e a Tabelas II.8 no Anexo II ).
A maioria ( 67,4 \% ) afirmou ter estudado e aprendido muitos ou todos os conteúdos.
Na análise por Grande Região, a proporção de respondentes que escolheram a opção não estudou ainda a maioria desses conteúdos foi pequena.
Observa -se que, nas regiões Norte ( 4,6 \% ), Nordeste ( 4,7 \% ) e Centro-Oeste ( 4,9 \% ), as proporções foram maiores do que a média nacional ( 4,2 \% ).
São observadas diferenças estatisticamente significativas entre as regiões Sudeste com as regiões Norte, Nordeste e Centro-Oeste, e entre a região Sul com as regiões Nordeste e Centro-Oeste.
Em todas as regiões, a maioria dos presentes afirmou ter estudado e aprendido muitos ou todos os conteúdos, com proporções variando entre 61,9 \%, na região Norte, a 70,1 \%, na região Sudeste.
Considerando -se, separadamente, as opiniões de estudantes em todos os quartos de desempenho, observa -se que, no quarto inferior, 7,8 \% ofereceram como resposta que não estudou ainda a maioria desses conteúdos, sendo 2,0 \% os do quarto superior com a mesma resposta, valores decrescentes com o aumento do desempenho.
As diferenças entre os alunos que optaram por esse motivo de dificuldade entre todos os quartos de desempenhos são estatisticamente significativas.
Tendo -se em conta o quarto superior, 79,4 \% dos alunos afirmaram ter estudado e aprendido muitos ou todos os conteúdos.
No outro extremo, no primeiro quarto, 54,4 \% optaram pelas mesmas categorias.
As proporções são crescentes com o aumento do desempenho. 4.7 TEMPO GASTO PARA CONCLUIR A PROVA Ao responderem sobre o tempo de conclusão da prova ( Questão 9 ), um pouco mais de três quartos dos estudantes ( 77,2 \% ) afirmaram ter gasto Entre duas e quatro horas ( Gráfico 4.17, Gráfico 4.18 e, no Anexo II, a Tabela II.9 ).
Considerando -se as cinco Grandes Regiões brasileiras, os percentuais dos que utilizaram Entre duas e quatro horas para finalizar a prova na região Sudeste ( 78,9 \% ) foi superior ao percentual nacional.
Nas demais Grandes Regiões, o percentual de alunos que dispensaram Entre duas e quatro horas para concluir a prova foi inferior ao percentual nacional, entre 75,4 \% e 77,1 \%, como mostrado no Gráfico 4.17.
Há diferença estatisticamente significativa entre as proporções da Grande Região Sudeste e as demais regiões.
Uma vez analisadas as alternativas escolhidas pelos estudantes que se situam nos diferentes quartos de desempenho.
Observa -se uma tendência crescente : uma proporção maior de participantes no quarto superior declarou ter gasto Entre duas e quatro horas para concluir a prova quando comparada com a proporção de participantes nos quartos inferiores.
Há diferenças estatisticamente significativas entre os quartos de desempenho.

% matched lemmas: 
\end{textsample}
